\chapter{Συμπεράσματα και Μελλοντική Εργασία}
\section{Σύνοψη της Ερευνητικής Προσπάθειας και Βασικά Ευρήματα}
Η παρούσα διπλωματική εργασία σχεδιάστηκε και υλοποιήθηκε με κεντρικό γνώμονα τη διερεύνηση, ανάπτυξη και πρακτική εφαρμογή μιας πλατφόρμας Ψηφιακού Διδύμου (\en{Digital Twin}) σε σχολικές υποδομές. Μέσα από τη δημιουργία της αναπτυχθείσας πλατφόρμας και την πιλοτική της εφαρμογή στα Εκπαιδευτήρια Πλάτων, αποδείχθηκε ότι η μετάβαση σε μια Έξυπνη Πανεπιστημιούπολη (\en{Smart Campus}) δεν αποτελεί απλώς μια θεωρητική ακαδημαϊκή έννοια, αλλά μια άμεσα υλοποιήσιμη, τεχνολογικά εφικτή και οικονομικά αποδοτική πραγματικότητα.

Τα λειτουργικά ευρήματα της έρευνας υποστήριξαν ισχυρά τις αρχικές ερευνητικές υποθέσεις. Στον άξονα της ενεργειακής διαχείρισης (\en{Energy Efficiency}), η ενσωμάτωση IoT μικροελεγκτών και η χρήση αυτοματισμών μέσω του \en{openHAB} οδήγησαν σε μείωση της μέσης ημερήσιας κατανάλωσης ηλεκτρικής ενέργειας κατά 19,7\% (από 223 kWh σε 179 kWh ανά ημέρα). Εντυπωσιακότερη υπήρξε η \en{Year-over-Year (YoY)} σύγκριση για τον μήνα Νοέμβριο, όπου η αποκοπή των κρυφών φορτίων (\en{phantom loads}) του κυλικείου απέδωσε εξοικονόμηση της τάξης του 32,5\%. Στον άξονα της Ποιότητας Εσωτερικού Περιβάλλοντος (\en{IEQ}), το σύστημα κλειστού βρόχου διατήρησε τα επίπεδα θορύβου των σχολικών αιθουσών κάτω από το κρίσιμο κατώφλι των 55 dBA, παρέχοντας μια μη-παρεμβατική οπτική ανατροφοδότηση (\en{quiet-lamp}) που ενίσχυσε τη μαθησιακή διαδικασία χωρίς να διασπά την προσοχή.

Σε τεχνοοικονομικό επίπεδο, το αναπτυχθέν μοντέλο επιβεβαίωσε την εμπορική βιωσιμότητα της λύσης. Με συνολικό χρόνο απόσβεσης (\en{Payback Period}) 42,5 μηνών και Πενταετή Απόδοση Επένδυσης (\en{5-year ROI}) 41,02\%, η πλατφόρμα κατέδειξε ότι η χρήση τεχνολογιών ανοιχτού κώδικα (\en{open-source}) και εξοπλισμού COTS (\en{Commercial Off-The-Shelf}) αποτελεί μια εξαιρετικά ανταγωνιστική εναλλακτική έναντι βιομηχανικών συστημάτων BMS με υψηλότερο αρχικό κόστος εγκατάστασης.

\section{Επιστημονική και Τεχνολογική Συνεισφορά}
Η σημαντικότερη συνεισφορά της παρούσας μελέτης εντοπίζεται στην επιτυχή γεφύρωση του χάσματος μεταξύ της απλής ψηφιακής τηλεμετρίας (\en{Digital Shadow}) και του αυθεντικού Ψηφιακού Διδύμου (\en{Digital Twin}). Η πλειονότητα των σύγχρονων εκπαιδευτικών IoT εφαρμογών περιορίζεται στη μονόδρομη οπτικοποίηση δεδομένων \cite{Kritzinger2018}. Η παρούσα πλατφόρμα, αντιθέτως, υλοποίησε μια αυστηρή αρχιτεκτονική πέντε επιπέδων και εγκαθίδρυσε αμφίδρομη ροή δεδομένων, κλείνοντας τον κύκλο \textit{\en{Sense-Decide-Act}}.

Μέσω της τεχνικής του 'Brownfield Integration', το σύστημα τεκμηρίωσε υψηλό βαθμό διαλειτουργικότητας (\en{interoperability}). Δεν αντικατέστησε τις βαριές υποδομές του σχολείου, αλλά ενοποιήθηκε με αυτές (όπως η σύνδεση με τα υφιστάμενα PLC των αυτοματοποιημένων πυλών οχημάτων), χρησιμοποιώντας το πρωτόκολλο MQTT ως κεντρικό δίαυλο επικοινωνίας. Το γεγονός αυτό, σε συνδυασμό με το προηγμένο 3D/2D γραφικό περιβάλλον, καθιστά το προτεινόμενο σύστημα ένα πρότυπο αρχιτεκτονικό μοντέλο, ικανό να επεκταθεί και να εφαρμοστεί σε οποιαδήποτε άλλη εκπαιδευτική ή κτιριακή εγκατάσταση.

\section{Περιορισμοί της Παρούσας Έρευνας}
Όπως κάθε εκτεταμένη τεχνολογική εγκατάσταση σε πραγματικό περιβάλλον λειτουργίας, η παρούσα μελέτη χαρακτηρίζεται από ορισμένους περιορισμούς, οι οποίοι ωστόσο παρέχουν πολύτιμα μαθήματα για μελλοντικές αναβαθμίσεις.

Αρχικά, η αρχιτεκτονική του φυσικού επιπέδου (\en{Edge Hardware}) βασίστηκε αποκλειστικά στο ασύρματο τοπικό δίκτυο (\en{Wi-Fi}) της σχολικής εγκατάστασης. Αυτή η σχεδιαστική επιλογή εξαρτά την αδιάλειπτη ροή της τηλεμετρίας (\en{uptime}) από τη σταθερότητα του εμπορικού δικτύου του κτιρίου (\en{Access Points, Routers}). Σε περιπτώσεις διακοπών δικτύου, υπήρξε προσωρινή απώλεια επικοινωνίας μεταξύ των κόμβων (\en{ESP32/Shelly}) και του κεντρικού διακομιστή (\en{MQTT Broker}).

Δεύτερον, ο αισθητηριακός εξοπλισμός χαμηλού κόστους, όπως τα μικρόφωνα I2S για τη μέτρηση της ηχητικής πίεσης (\en{dBA}), απαιτεί περιοδική βαθμονόμηση (\en{calibration}) έναντι πιστοποιημένων βιομηχανικών ηχομέτρων, προκειμένου να διατηρείται η ακρίβεια σε επίπεδα επιστημονικής αυστηρότητας. Τέλος, η αξιολόγηση πραγματοποιήθηκε σε μια μεμονωμένη σχολική εγκατάσταση. Μολονότι τα ευρήματα είναι πρακτικά σημαντικά, η γενίκευσή τους σε διαφορετικές κλιματικές ζώνες ή σε παλαιότερα, λιγότερο μονωμένα κτίρια, θα απαιτούσε την εφαρμογή του συστήματος σε ένα ευρύτερο δίκτυο σχολείων.

\section{Προτάσεις για Μελλοντική Έρευνα και Ανάπτυξη}
Η αρχιτεκτονική της αναπτυχθείσας πλατφόρμας σχεδιάστηκε να είναι αρθρωτή (\en{modular}) και επεκτάσιμη. Με βάση τα αποτελέσματα της παρούσας διπλωματικής εργασίας, αναδεικνύονται τρεις κεντρικοί άξονες για μελλοντική έρευνα και τεχνολογική επέκταση.

\subsection{Ενσωμάτωση Μοντέλων Πληροφοριών Κτιρίου (\en{BIM})}
Η τρέχουσα τρισδιάστατη οπτικοποίηση κατασκευάζεται διαδικαστικά (\en{procedurally}) μέσω τεχνολογιών web (\en{React Three Fiber}). Το επόμενο ερευνητικό βήμα είναι η πλήρης διασύνδεση του ψηφιακού διδύμου με Μοντέλα Πληροφοριών Κτιρίου (\en{Building Information Modeling - BIM}) και Γεωγραφικά Συστήματα Πληροφοριών (\en{GIS}). Η εισαγωγή προτύπων αρχείων, όπως το \en{Industry Foundation Classes (IFC)}, θα επιτρέψει στο ψηφιακό δίδυμο να κατέχει πλήρη κατασκευαστική επίγνωση των υλικών (π.χ. θερμοπερατότητα τοιχοποιίας, διατομές καλωδιώσεων), μετατρέποντάς το σε ένα ισχυρό εργαλείο Διαχείρισης Κύκλου Ζωής (\en{Lifecycle Management}) για το σχολικό κτίριο.

\subsection{Επέκταση Αισθητηριακού Δικτύου (\en{IEQ} \& \en{CO2})}
Στο πεδίο της Ποιότητας Εσωτερικού Περιβάλλοντος (\en{IEQ}), το σύστημα μπορεί να επεκταθεί πέραν της ακουστικής άνεσης. Η εγκατάσταση αισθητήρων ποιότητας αέρα (για μέτρηση Διοξειδίου του άνθρακα (\en{CO2}), Πτητικών Οργανικών Ενώσεων - \en{VOCs} και Αιωρούμενων Σωματιδίων - \en{PM2.5}) αποτελεί ερευνητική προτεραιότητα. Σε συνδυασμό με συστήματα Μηχανικού Αερισμού Ανάκτησης Θερμότητας (\en{HVAC/MVHR}), το ψηφιακό δίδυμο θα μπορεί να ελέγχει δυναμικά τον εξαερισμό της αίθουσας (\en{Demand-Controlled Ventilation}), συμβάλλοντας στη διατήρηση συνθηκών που συνδέονται με τη βελτίωση της γνωστικής απόδοσης των μαθητών.

\subsection{Εξέλιξη σε Γνωστικό Ψηφιακό Δίδυμο (\en{Cognitive Digital Twin}) μέσω Τεχνητής Νοημοσύνης (\en{AI})}
Ο τεράστιος όγκος των ιστορικών χρονοσειρών που συγκεντρώνονται πλέον στη βάση δεδομένων \en{InfluxDB} της πλατφόρμας, αποτελεί τον ιδανικό "καμβά" για την εφαρμογή αλγορίθμων Μηχανικής Μάθησης (\en{Machine Learning}). Το τελικό στάδιο εξέλιξης της πλατφόρμας είναι η μετάβαση από ένα ψηφιακό δίδυμο βασισμένο σε στατικούς κανόνες (\en{Rule-based DT}) σε ένα Γνωστικό Ψηφιακό Δίδυμο (\en{Cognitive Digital Twin}), με δυνατότητα σύνδεσης και σε σενάρια \en{XR}-βασισμένης εκπαίδευσης \cite{Li2024}.

Μέσω της Τεχνητής Νοημοσύνης, το σύστημα θα είναι σε θέση να εκτελεί Αναγνώριση Ανωμαλιών (\en{Anomaly Detection}) σε πραγματικό χρόνο. Για παράδειγμα, αναλύοντας τα μοτίβα κατανάλωσης ισχύος των συμπιεστών ψύξης, το σύστημα θα προβλέπει επικείμενες μηχανικές βλάβες πριν αυτές συμβούν (\en{Predictive Maintenance}), προτείνοντας αυτόματα εργασίες συντήρησης. Στο ίδιο πλαίσιο εντάσσεται και το \en{AI Energy Forecasting}, δηλαδή η πρόβλεψη μελλοντικών ενεργειακών απαιτήσεων και κόστους με χρήση ιστορικών χρονοσειρών και επιχειρησιακών μεταβλητών. Η εξέλιξη αυτή θα καταστήσει το ψηφιακό δίδυμο ικανό να λαμβάνει πλήρως αυτόνομες, προγνωστικές αποφάσεις, μεγιστοποιώντας τον χρόνο καλής λειτουργίας (\en{uptime}) των εγκαταστάσεων και προσφέροντας ένα πραγματικά "έξυπνο" τεχνολογικό οικοσύστημα για την εκπαίδευση του μέλλοντος.

\subsection{Ενσωμάτωση Τεχνολογιών Φωνητικών Βοηθών (\en{Voice Assistants})}
Μια ιδιαίτερα ενδιαφέρουσα κατεύθυνση μελλοντικής επέκτασης αφορά την ενσωμάτωση φωνητικών διεπαφών και τεχνολογιών επεξεργασίας φυσικής γλώσσας (\en{Natural Language Processing - NLP}) στο περιβάλλον του ψηφιακού διδύμου. Σε ένα τέτοιο σενάριο, ο διευθυντής ή ο διαχειριστής της σχολικής μονάδας δεν θα αλληλεπιδρά αποκλειστικά μέσω dashboards και μενού πλοήγησης, αλλά θα μπορεί να υποβάλλει φυσικά ερωτήματα ή εντολές, όπως «ποια είναι η κατανάλωση ενέργειας στο κυλικείο σήμερα;» ή «μείωσε τη θερμοκρασία στην αίθουσα 2». Αντί να απαιτείται πλοήγηση σε πίνακες ελέγχου, οι χρήστες θα μπορούν επίσης να θέτουν ερωτήματα σε καθημερινή γλώσσα, όπως «Ποια είναι η κατάσταση του συστήματος σήμερα;» ή «Υπάρχει κάποια ενεργειακή ανωμαλία;», με το σύστημα να παρέχει άμεσα ηχητική ή κειμενική σύνοψη της απόδοσης. Η προσέγγιση αυτή μπορεί να μειώσει σημαντικά το γνωστικό φορτίο χρήσης της πλατφόρμας και να καταστήσει το ψηφιακό δίδυμο περισσότερο προσιτό και φιλικό προς μη τεχνικούς χρήστες.

\subsection{Βελτιστοποίηση Υλικού και Λογισμικού (\en{Hardware \& Software Enhancements})}
Μια δεύτερη σημαντική κατεύθυνση αφορά τη συνεχή βελτιστοποίηση του ίδιου του τεχνολογικού υποβάθρου της πλατφόρμας. Σε επίπεδο υλικού, μπορούν να εξεταστούν πιο εξειδικευμένες edge υποδομές, όπως επιταχυντές τοπικής τεχνητής νοημοσύνης, οι οποίοι θα επιτρέπουν ταχύτερη προεπεξεργασία και εξαγωγή χαρακτηριστικών στο φυσικό πεδίο. Σε επίπεδο λογισμικού, η αξιοποίηση βελτιστοποιημένων μηχανισμών συγχρονισμού, όπως αποδοτικότερα κανάλια \en{WebSockets} και πιο προσαρμοσμένα real-time patterns, μπορεί να μειώσει περαιτέρω το latency, να αυξήσει τον ρυθμό δειγματοληψίας (\en{sampling rate}) και να βελτιώσει τη συνολική απόδοση (\en{performance}) του Ψηφιακού Διδύμου σε πραγματικό χρόνο.

\subsection{Αυτοματοποιημένη Προγνωστική Ανάλυση με Μηχανική Μάθηση (\en{Predictive Analytics with Machine Learning})}
Ο όγκος των τηλεμετρικών δεδομένων που έχει ήδη συγκεντρωθεί δημιουργεί ισχυρή βάση για την ανάπτυξη αυτοματοποιημένων μοντέλων προγνωστικής ανάλυσης. Πέρα από τις σχετικά απλές rule-based ειδοποιήσεις, μελλοντικά η πλατφόρμα μπορεί να αξιοποιεί αλγορίθμους Μηχανικής Μάθησης για την πρόβλεψη βλαβών και την εκτίμηση μελλοντικών ενεργειακών απαιτήσεων χωρίς χειροκίνητη ανάλυση από τον χρήστη. Ενδεικτικά, η συμπεριφορά του ψυγείου του κυλικείου, η ενεργειακή απόκριση των επιμέρους υποσυστημάτων και τα μοτίβα λειτουργίας ανά εποχή μπορούν να τροφοδοτήσουν μοντέλα \en{Predictive Analytics}, \en{Predictive Maintenance} και \en{AI Energy Forecasting}, ενισχύοντας σημαντικά τη δυνατότητα προληπτικής διαχείρισης της εγκατάστασης. Σε πιο προχωρημένο στάδιο, το \en{AI} θα μπορεί να αναλύει ιστορικά δεδομένα για να εντοπίζει μικρο-ανωμαλίες, όπως ανεπαίσθητους κραδασμούς στους συμπιεστές ή πτώσεις τάσης σε κρίσιμα φορτία, επιτρέποντας την έγκαιρη πρόβλεψη βλαβών και προστατεύοντας τον εξοπλισμό από κρίσιμες αστοχίες.

\subsection{Αυτόνομα Γνωστικά Ψηφιακά Δίδυμα (\en{Intelligent Cognitive Digital Twins})}
Η πλέον προωθημένη μελλοντική προοπτική αφορά τη μετάβαση από την υποβοήθηση αποφάσεων (\en{decision support}) στην ημιαυτόνομη ή πλήρως αυτόνομη λήψη αποφάσεων. Σε αυτό το εξελικτικό στάδιο, το ψηφιακό δίδυμο θα μπορεί να λειτουργεί ως \en{Intelligent Agent}, συνδυάζοντας τηλεμετρία, ιστορικά δεδομένα, προβλέψεις και επιχειρησιακούς κανόνες, ώστε να ενεργεί δυναμικά χωρίς άμεση ανθρώπινη παρέμβαση. Παραδείγματα τέτοιας λειτουργίας περιλαμβάνουν την αυτόματη ρύθμιση συστημάτων \en{HVAC} με βάση τον καιρό και την πληρότητα των χώρων, τη δυναμική ανακατανομή ενεργειακών φορτίων και την αυτενεργή βελτιστοποίηση συνθηκών λειτουργίας. Η κατεύθυνση αυτή συνδέεται άμεσα με το όραμα του \en{Cognitive Digital Twin}, δηλαδή ενός συστήματος που δεν αρκείται στην αναπαράσταση ή στη σύσταση ενεργειών, αλλά αναλαμβάνει ενεργά την προσαρμογή της ίδιας της υποδομής.

\section{Επίλογος}
Εν κατακλείδι, η παρούσα διπλωματική εργασία καταδεικνύει ότι ο ψηφιακός μετασχηματισμός των εκπαιδευτικών υποδομών δεν αποτελεί πλέον ένα θεωρητικό ή μακρινό όραμα, αλλά μια άμεσα υλοποιήσιμη τεχνολογική πραγματικότητα. Η σύζευξη τεχνολογιών αιχμής, όπως το Διαδίκτυο των Πραγμάτων (\en{IoT}) και τα Ψηφιακά Δίδυμα (\en{Digital Twins}), με προσεγγίσεις ανοιχτού κώδικα προσφέρει ένα ισχυρό οπλοστάσιο στα χέρια των σύγχρονων Μηχανικών. Η αναπτυχθείσα πλατφόρμα θέτει πρακτικά τα θεμέλια για τη δημιουργία Έξυπνων Πανεπιστημιουπόλεων (\en{Smart Campuses}) που δεν προστατεύουν απλώς την υγεία και τη μαθησιακή συγκέντρωση, αλλά ταυτόχρονα εξοικονομούν πόρους και σέβονται το περιβάλλον. Καθώς οι τεχνολογίες Τεχνητής Νοημοσύνης και Μηχανικής Μάθησης θα ενσωματώνονται ολοένα και βαθύτερα σε αυτά τα συστήματα τα επόμενα χρόνια, τα σχολεία του μέλλοντος αναμένεται να μετατραπούν σε ζωντανούς, δυναμικούς οργανισμούς, ικανούς να "αισθάνονται", να προσαρμόζονται και να αλληλεπιδρούν με τους μαθητές τους, διαμορφώνοντας ένα ασφαλέστερο, αποδοτικότερο και πιο βιώσιμο αύριο για την εκπαίδευση.
